\section{Ejercicios básicos con CLIPS}

\subsection{Ejercicio 2.1}

\textbf{Si cargamos la regla y el fichero de datos anterior, ¿cuántas activaciones se producen de la regla?}

Dos, una por cada hecho que encaja con \texttt{(EsPadre ?variable)}: f-1 (Pedro) y f-2 (Juan).

\begin{itemize}
\item \texttt{(run 1)}: se ejecuta solo la activación más reciente (Juan) y se añade \texttt{(QuiereASusHijos Juan)}. La de Pedro queda en la Agenda. \texttt{(reset)} restaura los hechos iniciales y vuelve a generar las dos activaciones.
\item \texttt{(run 2)}: se ejecutan las dos y se añaden \texttt{(QuiereASusHijos Juan)} y \texttt{(QuiereASusHijos Pedro)}.
\item \texttt{(run)}: se ejecutan todas las activaciones, que son estas dos, con el mismo resultado que \texttt{(run 2)}.
\end{itemize}

\subsection{Cuestiones 2.3}

Los hechos son \texttt{(progenitor X Y)} (X es padre o madre de Y) y \texttt{(hombre X)} o \texttt{(mujer X)}. Las relaciones deducidas se guardan como \texttt{(relacion <tipo> X Y)} (``X es <tipo> de Y''):

\begin{lstlisting}
(defrule hermanos
   (progenitor ?padre ?x) (hombre ?padre)
   (progenitor ?madre ?x) (mujer ?madre)
   (progenitor ?padre ?y&~?x) (progenitor ?madre ?y)
   => (assert (relacion hermano ?x ?y)))

(defrule tios
   (relacion hermano ?t ?p) (progenitor ?p ?s)
   => (assert (relacion tio ?t ?s)))

(defrule sobrinos
   (relacion tio ?t ?s)
   => (assert (relacion sobrino ?s ?t)))

(defrule primos
   (progenitor ?a ?x) (progenitor ?b ?y) (relacion hermano ?a ?b)
   => (assert (relacion primo ?x ?y)))

(defrule abuelos
   (progenitor ?a ?p) (progenitor ?p ?n)
   => (assert (relacion abuelo ?a ?n)))
\end{lstlisting}

Dos personas son hermanos cuando comparten madre y padre. Sobre esa relación se apoyan las demás, ya que un tío es el hermano de un progenitor, los primos son los hijos de dos hermanos y un abuelo es el progenitor de un progenitor.

\textbf{1. ¿Quiénes son los hermanos, tíos, sobrinos, primos y abuelos de Archie?}

Solo tiene abuelo, que es Carlos. No tiene hermanos, tíos, sobrinos ni primos. La regla de hermanos exige compartir padre y madre, y en los hechos no figura la madre de Guillermo y Enrique ni la de Archie y Lillibet.

\textbf{2. ¿Quiénes son abuelos y de quién?}

\begin{itemize}
\item Isabel y Felipe, de Guillermo y Enrique.
\item Carlos, de Jorge, Carlota, Luis, Archie y Lillibet.
\end{itemize}

\textbf{3. ¿Quiénes son tíos y de quién?}

Ana, Andrés y Eduardo son tíos de Guillermo y de Enrique.

\subsection{Cuestiones 2.4}

Se añaden los hechos nuevos:

\begin{lstlisting}
(mujer Camilla) (mujer Diana) (hombre Timothy) (mujer Sofia)
(progenitor Diana Guillermo) (progenitor Diana Enrique)
(pareja Carlos Camilla) (pareja Ana Timothy) (pareja Eduardo Sofia)
\end{lstlisting}

Con Diana como madre, Guillermo y Enrique pasan a ser hermanos sin necesidad de reglas nuevas. Las parejas sí requieren reglas nuevas, para el parentesco político:

\begin{lstlisting}
(defrule pareja-simetrica
   (pareja ?x ?y) => (assert (pareja ?y ?x)))

(defrule tios-politicos
   (pareja ?p ?t) (relacion tio ?t ?s)
   => (assert (relacion tio-politico ?p ?s)))

(defrule sobrinos-politicos
   (relacion tio-politico ?p ?s)
   => (assert (relacion sobrino-politico ?s ?p)))
\end{lstlisting}

Un tío político es la pareja de un tío, y el sobrino político es la relación simétrica. La regla de parejas hace falta porque los hechos solo dan una dirección de cada pareja.

\textbf{1. ¿Quiénes son los hermanos, tíos, sobrinos, primos y abuelos de Archie?}

\begin{itemize}
\item Hermanos: ninguno, porque sigue sin figurar su madre.
\item Tíos: Guillermo.
\item Sobrinos: ninguno.
\item Primos: Jorge, Carlota y Luis.
\item Abuelos: Carlos y Diana.
\end{itemize}

\textbf{2. ¿Quiénes son abuelos y de quién?}

\begin{itemize}
\item Isabel y Felipe, de Guillermo y Enrique.
\item Carlos y Diana, de Jorge, Carlota, Luis, Archie y Lillibet.
\end{itemize}

\textbf{3. ¿Quiénes son tíos y de quién?}

\begin{itemize}
\item Ana, Andrés y Eduardo, de Guillermo y Enrique.
\item Guillermo, de Archie y Lillibet.
\item Enrique, de Jorge, Carlota y Luis.
\item Tíos políticos: Timothy y Sofía, de Guillermo y Enrique.
\end{itemize}

Camilla no tiene ninguna de estas relaciones, porque no es progenitora de los hijos de Carlos y Carlos no es tío de nadie.
